\section{Related Work}
% ──────────────────────────────────────────────────────────────────────────────

Efforts to taxonomise XAI methods have produced a substantial but fragmented body of work. Arrieta et al.~\cite{Arrieta2020XAI} established the foundational distinction between interpretable-by-design models and post-hoc explanations. Schwalbe and Finzel~\cite{schwalbe2024comprehensive} synthesised over fifty surveys into a framework showing that properties such as faithfulness and stability are widely referenced but inconsistently operationalised. Speith~\cite{speith2022review} identified a structural source of this inconsistency: taxonomies grounded in mechanisms versus outputs classify the same methods differently, limiting cross-study comparability.

Medical XAI taxonomies address this gap but remain incomplete. Jin et al.~\cite{jin2023guidelines} distinguish technical faithfulness from clinical usefulness; Holzinger et al.~\cite{holzinger2022information} emphasise robustness under distributional shift as unresolved for clinical deployment; Tjoa and Guan~\cite{Tjoa2021MedicalXAI} survey XAI methods broadly without formalising evaluation criteria. Varghese et al.~\cite{Varghese2025OnDE} rank XAI metrics using Borda count and multiple correlation techniques, identifying faithfulness and sensitivity as primary evaluative axes without cross-domain validation. Reviews in radiology and nuclear medicine~\cite{groen2022systematic,de2023explainable,saw2025current} further demonstrate that evaluation metrics are inconsistently reported and rarely subjected to comparative analysis. Reyes et al.~\cite{reyes2020interpretability} argue that explanations must align with clinical reasoning---a requirement frequently absent from benchmark-driven evaluations.

The metrics literature exposes practical consequences of this fragmentation. Rosenfeld~\cite{rosenfeld2021better} demonstrates that most evaluations prioritise algorithmic fidelity while neglecting human-centred dimensions. Kadir et al.~\cite{kadir2023evaluation} confirm that sensitivity, faithfulness, and robustness dominate, with limited attention to human factors. Phuong et al.~\cite{phuong2023benchmarking} and Klein et al.~\cite{Klein2024NavigatingTM} provide evidence that without structured taxonomies, metric selection remains ad hoc. Proposed frameworks such as Mirzaei et al.~\cite{mirzaei2023xai} offer selection criteria but lack domain-specific grounding. Complementary analyses~\cite{coroama2022evaluation,boggust2023saliency} reinforce that the absence of shared evaluation vocabulary impedes meaningful synthesis.

Human-centred evaluation has received increasing attention. Chen et al.~\cite{chen2022explainable} identify the systematic absence of human-centred components in medical XAI pipelines. Bauer and Michalowski~\cite{bauer2025human} highlight cognitive load, trust calibration, and decision accuracy as key constructs requiring operationalisation. Dey et al.~\cite{dey2022human} document continued methodological fragmentation, and M{\"u}ller et al.~\cite{muller2022explainability} emphasise that regulatory demands for explainability and causability extend beyond standard faithfulness metrics. Recent surveys~\cite{houssein2025explainable,zhang2026comprehensive,Ahmed2026XAI} confirm that standardised evaluation frameworks remain underdeveloped, with Ahmed et al.~\cite{Ahmed2026XAI} identifying metric standardisation as a major open challenge.

What remains absent is a unified taxonomy grounded in empirical usage patterns rather than purely conceptual design. No existing framework jointly spans oncology, cardiology, neuroimaging, radiology, and clinical decision support while integrating mathematical faithfulness, system robustness, and human-centred trust. The present study addresses this gap through systematic co-occurrence analysis and empirically validated taxonomy construction.
